\section{Verification purpose}

This case verifies the coupling of wind-adjusted surface spread to the Huygens and level-set propagation geometry.  Uniform fuels and flat terrain isolate wind as the only directional forcing.  The verification question is whether the simulated head rate is within 10 percent of \(3.635~\mathrm{m\,min^{-1}}\) and the unclipped arrival-time isochrone has length-to-width ratio within 10 percent of 1.24.

The scientific setup and historical expected response are reconstructed from Section~3.2 of the supplied \emph{ELMFIRE Guide}. The Guide is used as reference provenance, not as an instruction or a substitute for evidence. Configuration statements below are checked against the generated GeoTIFFs and namelist; current results and decisions come only from case-specific outputs and generated metrics. Where the generated experiment differs from the historical example, the report states that difference explicitly.

\section{Mathematical formulation and numerical implementation}

The normal-velocity construction combines the backing, flanking, and heading rates into an expanding ellipse.  Its major axis aligns with the wind-driven heading direction, while the ratio \(L/W\) characterizes anisotropy independent of absolute size.  ELMFIRE writes rate in metres per minute and arrival time on the same raster grid used to infer the ellipse.

For uncovered fuel of depth \(H\) ft, the Guide uses
\begin{equation}
WAF_u=\frac{1.83}{\ln[(20+0.36H)/(0.13H)]}, \qquad U_{mf}=WAF_uU_{20}.
\end{equation}
The resulting effective midflame wind determines the ellipse aspect ratio
\begin{equation}
\frac{L}{W}=\min\left(0.936e^{0.1147U_{mf,e}}+0.461e^{-0.0692U_{mf,e}}-0.397,8\right).
\end{equation}
The local head rate is \(R_H=R_0(1+\phi_w)\); the Huygens/level-set construction uses the ellipse to obtain normal speed away from the heading direction.

\section{Derivation of expected behavior}

Constant forcing produces self-similar elliptical isochrones: scaling time changes both axes proportionally and leaves \(L/W\) constant.  The configured Fuel Model 5 response supplies the independent head-rate target \(3.635~\mathrm{m\,min^{-1}}\), and the GUIDE construction supplies \(L/W=1.24\).

% BEGIN EXPLICIT METRIC DERIVATION
For the generated Fuel Model~5 state, the independent no-wind rate is
\(R_0=0.41042~\mathrm{m\,min^{-1}}\).  The \(6~\mathrm{mph}\) 20-ft wind
and 2-ft fuel-bed depth give a wind-adjustment factor
\(WAF=0.418640\), hence
\(U_{mf}=6(0.418640)=2.51184~\mathrm{mph}\).  With the calculated wind
factor \(\phi_w=7.87958\),
\[
 R_h=R_0(1+\phi_w)=0.41042(1+7.87958)
     =3.6444~\mathrm{m\,min^{-1}},
\]
consistent with the Guide reference \(3.635~\mathrm{m\,min^{-1}}\).
The ellipse length-to-width reference uses
\[
 \frac{L}{W}=0.936e^{0.2566U}+0.461e^{-0.1548U}-0.397,
\]
where \(U\) is in \(\mathrm{m\,s^{-1}}\).  Here
\(U=2.51184~\mathrm{mph}(0.44704)=1.12291~\mathrm{m\,s^{-1}}\), so
substitution gives \(L/W=1.23898\), reported as 1.24.  The report evaluates the relative substitutions
\(e_R=|R_h-3.635|/3.635\) and
\(e_{LW}=|(L/W)-1.24|/1.24\) against their tabulated limits.
% END EXPLICIT METRIC DERIVATION

\section{Verification metrics and acceptance criteria}

Maximum finite spread rate measures the heading response.  The shape metric is \(\sqrt{\lambda_{\max}/\lambda_{\min}}\), where the eigenvalues come from the covariance of burned-cell coordinates at the last isochrone that remains away from the boundary.  This metric is rotation invariant and avoids a wind-axis assumption.  Each component uses 10 percent relative tolerance, and both must pass.

\section{Simulation configuration and predicted outcome}

The \(126\times126\), 30 m domain uses uniform Fuel Model 5, 6 mph wind, flat terrain, and an upwind point ignition.  The predicted arrival field is an ellipse elongated downwind, with head rate of spread (ROS) \(3.635~\mathrm{m\,min^{-1}}\) and \(L/W=1.24\).

\begin{center}
\begin{tabular}{@{}p{0.20\textwidth}p{0.20\textwidth}p{0.12\textwidth}p{0.38\textwidth}@{}}
\toprule
Parameter & Configured value & Units & Role in the experiment \\
\midrule
Grid & \(126\times126\) at 30 & m per cell & Resolves the whole spatial response on the generated north-up raster \\
Projection & EPSG:32635 & --- & Fixes the map-coordinate and row-orientation convention \\
Time control & \(\Delta t=1\), \(t_{stop}=259200\) & s & Whole-second step; stop time is an exact timestep multiple \\
Meteorology interval & 3600 & s & Defines weather-band timing and stalled-final compatibility \\
Live moisture & herb 60; woody 90 & percent & Holds live-fuel damping fixed while the named mechanism is isolated \\
Evidence & prepared inputs; final case outputs & --- & Separates configuration provenance from acceptance evidence \\
\bottomrule
\end{tabular}
\end{center}

Figure~\ref{fig:input-configuration} shows the prescribed spatial condition used in this experiment. The displayed field represents the initial level-set geometry or a controlling fuel, terrain, moisture, wind, or canopy quantity, as identified by the figure labels. The plotted values come from the prepared simulation inputs, not from the expected response.

\begin{figure}[H]
  \centering
  \IfFileExists{../figures/input_configuration.pdf}{%
    \includegraphics[width=0.98\textwidth,height=0.46\textheight,keepaspectratio]{input_configuration.pdf}%
  }{\fbox{\parbox{0.88\textwidth}{Prepared-input evidence is unavailable. The prescribed configuration is described in the text; this absence does not establish that it was simulated.}}}
  \caption{Prepared input conditions from the prescribed spatial fields. Coordinates are in metres; the colorbar identifies the displayed field and its physical units or category. The ignition panel is derived directly from \(\phi\leq0\).}
  \label{fig:input-configuration}
\end{figure}

\section{Actual simulation results}

The required evidence is generated only from case-specific ELMFIRE outputs. The current postprocessing status is \textbf{\VerificationStatus}; \PassingMetricCount{} of \MetricCount{} required metric rows pass across \VariantCount{} simulation condition(s). A missing output is reported as NOT RUN or NOT EVALUABLE and is not interpreted as verification evidence.

\begin{figure}[H]
  \centering
  \IfFileExists{../figures/whole_domain_result.pdf}{%
    \includegraphics[width=0.82\textwidth,height=0.58\textheight,keepaspectratio]{whole_domain_result.pdf}%
  }{\fbox{\parbox{0.85\textwidth}{No case-specific whole-domain ELMFIRE raster is available.}}}
  \caption{Final whole-domain ELMFIRE time of arrival for the displayed simulation condition, in seconds. The coordinate reference system is stated in the figure. The calculated metrics, not visual appearance, determine acceptance.}
\end{figure}

\section{Calculated metrics and verification decision}

\begin{center}
\begin{tabular}{@{}llll@{}}
\toprule
Simulation condition / metric & Expected & Calculated & Status \\
\midrule
\MetricRows
\bottomrule
\end{tabular}
\end{center}

The displayed overall decision is \textbf{\VerificationStatus}. When calculated verification and validation metrics exists, result analysis procedure generates every table value and status; when it is absent, the report states NOT EVALUABLE with no invented rows.

\section{Technical interpretation}

A rate-only discrepancy points toward wind adjustment or unit conversion; a shape-only discrepancy points toward wavelet geometry, arrival-time extraction, or boundary clipping. No scientific conclusion should be drawn when the status is NOT RUN or NOT EVALUABLE.

\section{Scope and limitations}

The test covers one steady wind speed and one fuel model. The covariance estimator is intentionally applied only before boundary contact. This case provides implementation verification evidence only.
